% GENERATED FILE. Edit canonical lessons, then run npm run generate:books.
% canonical-source-hash: fnv1a64:5914b09b63ab64c2
% canonical-lessons: HI-C420-dhalan, HI-C420-galiyara, HI-C420-ghusna, HI-C420-lurhakna, HI-C420-garnaintr

\chapter{Slopes and Passages}
\label{ch:hi-slopes-and-passages}
\hlchaptermodality{420}

\begin{chapteropening}
\textbf{By the end of this chapter:} \emph{I can go up and down a slope, push through a passage, roll downhill and say something is digging into my foot.}

Complete the last lesson of chapter 420: I can go up and down a slope, push through a passage, roll downhill and say something is digging into my foot.
\end{chapteropening}

\section[\texorpdfstring{\hi{ढलान} (ḍhalān)}{ढलान (ḍhalān)}]{\hi{ढलान} (ḍhalān) — a slope}

\label{lesson:HI-C420-dhalan}



\begin{quote}
Before the new one: say the Hindi for a building, then the Hindi for a signboard.
\end{quote}

\subsection*{You'll want to know: \hi{ढलान}}
\textbf{\hi{ढलान}} — \emph{ḍhalān} — \textquotedblleft{}a slope\textquotedblright{}.

\begin{quote}
\hi{ढलान} \hi{पर} \hi{गाड़ी} \hi{धीरे} \hi{चलाइए।}
\end{quote}

\begin{quote}
\emph{ḍhalān par gāṛī dhīre calāiye.}
\end{quote}

\begin{quote}
Drive slowly on the slope.
\end{quote}

Feminine.

\begin{cousinweb}
\textbf{\hi{ढलान}} (\emph{ḍhalān}) is built on the stem of \hi{ढलना}, \textquotedblleft{}to slope, to go down\textquotedblright{}, which goes back to \emph{*ḍhalati}, \textquotedblleft{}bends over, falls\textquotedblright{}, a form CDIAL reconstructs (CDIAL 5581); CDIAL does not list the noun itself.
\end{cousinweb}

\begin{grammarlens}[title={What you've built}]
One more word for getting around.
\end{grammarlens}

\subsection*{Your turn}
\begin{itemize}
\raggedright
  \item \emph{Say it:} \emph{ḍhalān}
  \item \emph{Say it:} \emph{ḍhalān}, once more
  \item \emph{Say it:} \emph{ḍhalān par gāṛī dhīre calāiye.} — Drive slowly on the slope.
  \item \emph{Recall:} say the Hindi for a building, then the Hindi for a signboard, then say \emph{ḍhalān} again
\end{itemize}

\begin{tcolorbox}[breakable,colback=teal!4,colframe=teal!35!black,title={Before you move on}]
What does \hi{ढलान} mean? (\textquotedblleft{}A slope\textquotedblright{}.) Say it once more.
\end{tcolorbox}
\section[\texorpdfstring{\hi{गलियारा} (galiyārā)}{गलियारा (galiyārā)}]{\hi{गलियारा} (galiyārā) — a corridor; a passage}

\label{lesson:HI-C420-galiyara}



\begin{quote}
Before the new one: say the Hindi for a signboard, then the Hindi for a slope.
\end{quote}

\subsection*{You'll want to know: \hi{गलियारा}}
\textbf{\hi{गलियारा}} — \emph{galiyārā} — \textquotedblleft{}a corridor; a passage\textquotedblright{}.

\begin{quote}
\hi{ट्रेन} \hi{के} \hi{गलियारे} \hi{में} \hi{लोग} \hi{खड़े} \hi{थे।}
\end{quote}

\begin{quote}
\emph{ṭren ke galiyāre meṁ log khaṛe the.}
\end{quote}

\begin{quote}
People were standing in the corridor of the train.
\end{quote}

Masculine.

\begin{cousinweb}
\textbf{\hi{गलियारा}} (\emph{galiyārā}) is built on \hi{गली}, \textquotedblleft{}a lane\textquotedblright{}, from \emph{*galī}, \textquotedblleft{}a defile, a lane\textquotedblright{}, a form CDIAL reconstructs (CDIAL 4085): a lane inside a building or a train. CDIAL does not list galiyārā itself.
\end{cousinweb}

\begin{grammarlens}[title={What you've built}]
One more word for getting around.
\end{grammarlens}

\subsection*{Your turn}
\begin{itemize}
\raggedright
  \item \emph{Say it:} \emph{galiyārā}
  \item \emph{Say it:} \emph{galiyārā}, once more
  \item \emph{Say it:} \emph{ṭren ke galiyāre meṁ log khaṛe the.} — People were standing in the corridor of the train.
  \item \emph{Recall:} say the Hindi for a signboard, then the Hindi for a slope, then say \emph{galiyārā} again
\end{itemize}

\begin{tcolorbox}[breakable,colback=teal!4,colframe=teal!35!black,title={Before you move on}]
What does \hi{गलियारा} mean? (\textquotedblleft{}A corridor; a passage\textquotedblright{}.) Say it once more.
\end{tcolorbox}
\section[\texorpdfstring{\hi{घुसना} (ghusnā)}{घुसना (ghusnā)}]{\hi{घुसना} (ghusnā) — to push in; to get in}

\label{lesson:HI-C420-ghusna}



\begin{quote}
Before the new one: say the Hindi for a slope, then the Hindi for a corridor.
\end{quote}

\subsection*{You'll want to know: \hi{घुसना}}
\textbf{\hi{घुसना}} — \emph{ghusnā} — \textquotedblleft{}to push in; to get in\textquotedblright{}.

\begin{quote}
\hi{भीड़} \hi{में} \hi{घुसकर} \hi{हमने} \hi{बस} \hi{पकड़ी।}
\end{quote}

\begin{quote}
\emph{bhīṛ meṁ ghuskar hamne bas pakṛī.}
\end{quote}

\begin{quote}
We pushed through the crowd and caught the bus.
\end{quote}

Past (he …): \hi{वह} \hi{कमरे} \hi{में} \hi{घुसा}, \emph{vah kamre meṁ ghusā}.

\begin{cousinweb}
\textbf{\hi{घुसना}} (\emph{ghusnā}) goes back to \emph{*ghuss}, \textquotedblleft{}thrust in, pierce\textquotedblright{}, a form CDIAL reconstructs (CDIAL 4492), which lists Hindi \emph{ghusnā}, \textquotedblleft{}to be thrust in\textquotedblright{}.
\end{cousinweb}

\begin{grammarlens}[title={What you've built}]
One more word for getting around.
\end{grammarlens}

\subsection*{Your turn}
\begin{itemize}
\raggedright
  \item \emph{Say it:} \emph{ghusnā}
  \item \emph{Say it:} \emph{ghusnā}, once more
  \item \emph{Say it:} \emph{bhīṛ meṁ ghuskar hamne bas pakṛī.} — We pushed through the crowd and caught the bus.
  \item \emph{Recall:} say the Hindi for a slope, then the Hindi for a corridor, then say \emph{ghusnā} again
\end{itemize}

\begin{tcolorbox}[breakable,colback=teal!4,colframe=teal!35!black,title={Before you move on}]
What does \hi{घुसना} mean? (\textquotedblleft{}To push in; to get in\textquotedblright{}.) Say it once more.
\end{tcolorbox}
\section[\texorpdfstring{\hi{लुढ़कना} (luṛhaknā)}{लुढ़कना (luṛhaknā)}]{\hi{लुढ़कना} (luṛhaknā) — to roll (down)}

\label{lesson:HI-C420-lurhakna}



\begin{quote}
Before the new one: say the Hindi for a corridor, then the Hindi for to push in.
\end{quote}

\subsection*{You'll want to know: \hi{लुढ़कना}}
\textbf{\hi{लुढ़कना}} — \emph{luṛhaknā} — \textquotedblleft{}to roll (down)\textquotedblright{}.

\begin{quote}
\hi{पत्थर} \hi{ढलान} \hi{से} \hi{नीचे} \hi{लुढ़क} \hi{गया।}
\end{quote}

\begin{quote}
\emph{patthar ḍhalān se nīce luṛhak gayā.}
\end{quote}

\begin{quote}
The stone rolled down the slope.
\end{quote}

Past (it …): \hi{वह} \hi{गेंद} \hi{सड़क} \hi{पर} \hi{लुढ़की}, \emph{vah gẽd saṛak par luṛhkī}.

\begin{cousinweb}
\textbf{\hi{लुढ़कना}} (\emph{luṛhaknā}) is traced to Sanskrit \emph{luṭhati}, \textquotedblleft{}rolls, wallows\textquotedblright{} (CDIAL 11079), which lists Hindi \emph{luṛhaknā}, \textquotedblleft{}to roll, slip\textquotedblright{}.
\end{cousinweb}

\begin{grammarlens}[title={What you've built}]
One more word for getting around.
\end{grammarlens}

\subsection*{Your turn}
\begin{itemize}
\raggedright
  \item \emph{Say it:} \emph{luṛhaknā}
  \item \emph{Say it:} \emph{luṛhaknā}, once more
  \item \emph{Say it:} \emph{patthar ḍhalān se nīce luṛhak gayā.} — The stone rolled down the slope.
  \item \emph{Recall:} say the Hindi for a corridor, then the Hindi for to push in, then say \emph{luṛhaknā} again
\end{itemize}

\begin{tcolorbox}[breakable,colback=teal!4,colframe=teal!35!black,title={Before you move on}]
What does \hi{लुढ़कना} mean? (\textquotedblleft{}To roll (down)\textquotedblright{}.) Say it once more.
\end{tcolorbox}
\section[\texorpdfstring{\hi{गड़ना} (gaṛnā)}{गड़ना (gaṛnā)}]{\hi{गड़ना} (gaṛnā) — to dig into; to sink in (of something pointed)}

\label{lesson:HI-C420-garnaintr}



\begin{quote}
Before the new one: say the Hindi for to push in, then the Hindi for to roll.
\end{quote}

\subsection*{You'll want to know: \hi{गड़ना}}
\textbf{\hi{गड़ना}} — \emph{gaṛnā} — \textquotedblleft{}to dig into; to sink in (of something pointed)\textquotedblright{}.

\begin{quote}
\hi{जूते} \hi{की} \hi{कील} \hi{मेरे} \hi{पैर} \hi{में} \hi{गड़} \hi{रही} \hi{थी।}
\end{quote}

\begin{quote}
\emph{jūte kī kīl mere pair meṁ gaṛ rahī thī.}
\end{quote}

\begin{quote}
The nail in my shoe was digging into my foot.
\end{quote}

Past (it …): \hi{वह} \hi{काँटा} \hi{पैर} \hi{में} \hi{गड़} \hi{गया}, \emph{vah kā̃ṭā pair meṁ gaṛ gayā}.

\begin{cousinweb}
\textbf{\hi{गड़ना}} (\emph{gaṛnā}) is filed by CDIAL under Sanskrit \emph{gaḍati}, \textquotedblleft{}drops, runs as liquid\textquotedblright{} (CDIAL 3973), which lists Hindi \emph{gaṛnā}, \textquotedblleft{}to sink into\textquotedblright{}; CDIAL does not say how the one sense led to the other, so take the link as probable.
\end{cousinweb}

\begin{grammarlens}[title={What you've built}]
One more word for getting around.
\end{grammarlens}

\subsection*{Your turn}
\begin{itemize}
\raggedright
  \item \emph{Say it:} \emph{gaṛnā}
  \item \emph{Say it:} \emph{gaṛnā}, once more
  \item \emph{Say it:} \emph{jūte kī kīl mere pair meṁ gaṛ rahī thī.} — The nail in my shoe was digging into my foot.
  \item \emph{Recall:} say the Hindi for to push in, then the Hindi for to roll, then say \emph{gaṛnā} again
\end{itemize}

\begin{tcolorbox}[breakable,colback=teal!4,colframe=teal!35!black,title={Before you move on}]
What does \hi{गड़ना} mean? (\textquotedblleft{}To dig into; to sink in (of something pointed)\textquotedblright{}.) Say it once more.
\end{tcolorbox}
